\clearpage\fancyhead[L]{\small PALLAS.}\SourceLeaf{354}
\begin{center}\textbf{PALLAS.}\par\vspace{4pt}\hrule\vspace{4pt}\fontsize{10}{12}\selectfont Monatliche Correspondenz. Band XXV. Seite 389. 1812. April.\par\vspace{4pt}\hrule\vspace{4pt}\end{center}

\begin{flushright}Göttingen, den 15. April 1812.\end{flushright}
\par\indent Mit grosser Ungeduld erwartete ich dies Jahr die Wiedererscheinung der Pallas, um zu sehen, ob die Störungsrechnungen sich diesmal eben so schön bestätigen würden, wie im vorigen Jahre. Mondschein und schlechtes Wetter erlaubten mir aber nicht, früher sie zu sehen, als am 4. April; auch konnte ich sie an diesem nur kurze Zeit heitern Abend blos sehen, und erst am 9. April gelang mir eine Beobachtung, bei der aber die Declination blos geschätzt ist.

\begin{center}\fontsize{10}{12}\selectfont\setlength{\tabcolsep}{0pt}\renewcommand{\arraystretch}{1.15}\begin{tabular}{l}
1812 April 9\quad $11^{\mathrm u}\;32^{\mathrm m}\;27^{\mathrm s}$ AR. $\PallasSymbol\;269^\circ\;6'\;6''7$.\quad Decl. $16^\circ\;38'$
\end{tabular}\end{center}

\par\indent Diese Beobachtung zeigt wenigstens, dass die von Hrn. \textsc{Nicolai} berechnete Ephemeride vollkommen zutrifft. Es ist sehr zu wünschen, dass man die am 10. Juni einfallende Opposition nicht versäumen möge: allein die Beobachtung mit fixen Instrumenten wird wegen der Lichtschwäche des Planeten nicht leicht sein. Jetzt hat er nur die 11. Grösse.

\par\indent Die erste Berechnung der periodischen und Säcularstörungen durch Jupiter habe ich jetzt ganz vollendet, die Anzahl aller entwickelten Gleichungen steigt nahe an 400, wovon die grösste von $5\,\JupiterSymbol-2\,\PallasSymbol$ abhängende ungefähr 56 Min. beträgt. Es ist eine Arbeit von mehrern Tagen, nur einen einzigen Ort der Pallas zu berechnen. Von den Resultaten dieser Arbeit sind einige von hohem Interesse.

\par\indent Von der gestern vorgefallenen Bedeckung $\alpha$ Tauri konnte wegen bedeckten Himmels hier nichts gesehen werden.
